\section{Experimental design}
\label{sec:design}

\subsection{Conditions}

The campaign is the cross of two access levels and four agent frameworks:
$\{\Ltwo{}, \Lthree{}\} \times \{\hn{ursa}, \hn{dspy}, \hn{pi}, \hn{cursor}\}$,
eight cells, replicated --- plus a budget-matched \Lzero{} reference arm. We
attempted 59 runs and scored 49.

Replicate counts are deliberately unequal: two substrates at $n=10$ and two at
$n=5$. The campaign was budgeted in dollars rather than in runs, and one
substrate costs an order of magnitude more per run than another
(Section~\ref{sec:cost}). Failures reduce some cells further. Every denominator
in this paper counts \emph{attempts}, not survivors, which is the conservative
choice and is stated wherever a rate appears.

\subsection{Holding the model fixed}

Every cell runs the same language model, \code{gpt-5.2}. This is the point of
the design. In an earlier campaign of our own, one harness ran a different
model family from the others, and the resulting comparison was uninterpretable:
any difference could have been the framework or the model, and nothing in the
data could separate them.

Pinning it is fiddlier than it sounds, because the four adapters disagree about
how a model is named --- two take \code{openai:gpt-5.2}, one takes
\code{openai/gpt-5.2}, and the command-line agent takes a bare \code{gpt-5.2}.
We pin the name per adapter inside the frozen task specification rather than
passing a single command-line override, so that the model string is recorded
three independent times per run and the control is auditable after the fact. A
fifth adapter we maintain is excluded from the campaign entirely for the same
reason: it cannot run the pinned model, and including it would have reintroduced
exactly the confound the design exists to remove.

Framework versions were \hn{ursa} 0.15.1, \hn{dspy} 3.2.1, \hn{pi} 0.73.1 and
\hn{cursor} 2026.09.18, against Open FUSION Toolkit 26.6. Those are recorded
from the machine, not from the run records --- which stamp the task hash, the
prompt version, the model and the replicate index, but not the substrate's own
version. That is a gap in our harness and we would close it in a future
campaign.

\subsection{Equal compute budgets, and why they are hard to enforce}
\label{sec:budgets}

Giving every substrate the same budget turned out to be the hardest engineering
requirement in the benchmark, and the failure modes generalise well past this
paper.

Two adapters accepted a \code{timeout\_seconds} argument and silently ignored
it. Fixing that exposed two deeper problems. First, \textbf{a timeout raised as
an ordinary exception is swallowed}: agent frameworks wrap their step loops in
broad \code{except Exception} handlers and treat anything caught as a
recoverable step failure, so a timeout delivered by signal vanishes into the
agent's own retry logic --- and because \code{signal.alarm} is one-shot, a
single swallowed alarm disables the cap permanently. One harness absorbed its
90-minute cap this way and ran for over ten hours. Second, \textbf{killing the
agent process does not stop the work}: Python's \code{subprocess.run}, called
with a timeout, kills the direct child and then waits on its pipes again --- but
the solver's grandchildren inherit those pipes and hold them open, so the
timeout path itself blocks. One cell sat in that deadlock for 9 hours 19
minutes.

The fixes are to raise a \code{BaseException} subclass, for the same reason
\code{KeyboardInterrupt} is one, and to run each agent in its own process group
and signal the group. The general lesson is worth stating plainly: \textbf{an
agent harness is third-party code and cannot be trusted to honour a deadline
from inside}. We enforce the budget at three independent levels --- in-adapter,
by process-group kill, and by a driver-side watchdog that nothing running inside
the cell can intercept.

We also equalised the outer feedback-round setting to 1. Only two of the four
adapters implement outer review-and-fix rounds at all, so any higher value hands
those two a second pass the others never get --- an arbitrary two-tier design
that was the mechanical cause of their longer runtimes and larger spend in
earlier campaigns. One is the only value all four honour identically.

\subsection{What is held fixed}

The prompt text (version 3, hashed per run), the evaluation grid, the
60-parameter test set, the gate set and the scoring rule. The \Ltwo{} and
\Lthree{} problem statements are byte-identical except for ten passages, every
one of which concerns access level; they are reproduced with the differences
marked in Appendix~\ref{app:tasktext}. Results are pooled only within one prompt
version and one scoring epoch, and the aggregation tool refuses to mix them.

\subsection{Analysis, and one change we made after seeing the data}
\label{sec:analysis}

The unit of analysis is the cell. Pass rates carry Wilson score
intervals~\citep{wilson1927probable}, which --- unlike the normal
approximation --- do not collapse to zero width at $0/n$ or $n/n$.
Relative-$L_2$ is summarised by medians with percentile bootstrap
intervals~\citep{efron1979bootstrap} rather than means, because a single
catastrophic replicate makes a mean uninterpretable and we have several.

The \Ltwo{}/\Lthree{} contrast is paired within harness, since the harness is a
blocking factor rather than a treatment. We pre-specified an exact two-sided
sign test over the four harness pairs.

That test cannot return a $p$ below $0.125$ at four pairs, whatever the data
show, and the campaign produced a unanimous $4/4$ direction. We therefore
\textbf{changed the primary test after seeing the data}, to a stratified rank
statistic in the van Elteren family~\citep{vanelteren1960combination}, which
uses all 49 scored runs while still permuting labels only within harness. This
is a post-hoc change and we treat it as one: the analysis tool prints all three
tests on every run, all three are reported in Section~\ref{sec:l2l3}, and the
conclusion carries the caveat. Per-cell orderings are exploratory throughout;
the pooled level contrast is the claim that carries weight.
